\documentclass[10pt,aspectratio=169]{beamer}
% Preámbulo modular para presentaciones Beamer (Modelación Estadística)
% Diseño y paleta institucional del Tecnológico de Monterrey

\documentclass[aspectratio=169,xcolor={dvipsnames,table}]{beamer}

\usepackage[utf8]{inputenc}
\usepackage[spanish,mexico]{babel}
\usepackage{amsmath,amssymb,amsfonts,mathtools}
\usepackage{graphicx}
\usepackage{listings}
\usepackage{booktabs}
\usepackage{tikz}
\usetikzlibrary{matrix,arrows,decorations.pathmorphing,shapes.geometric,positioning}

% Paleta Institucional Tecnológico de Monterrey
\definecolor{TecAzul}{HTML}{0039A6}       % Azul TEC: color primario institucional
\definecolor{TecAzulOscuro}{HTML}{1A2E51} % Azul marino oscuro
\definecolor{TecRojo}{HTML}{EC2661}       % Rojo de acento (paleta miTec)
\definecolor{TecGris}{HTML}{646464}       % Gris neutro para comentarios y subtítulos
\definecolor{TecFondoBloque}{HTML}{F4F6F9}% Fondo suave para bloques

% Configuración de Tema Beamer
\usetheme{Boadilla}
\usecolortheme[named=TecAzulOscuro]{structure}
\setbeamercolor{alerted text}{fg=TecRojo}
\setbeamercolor{palette primary}{bg=TecAzulOscuro,fg=white}
\setbeamercolor{palette secondary}{bg=TecAzul,fg=white}
\setbeamercolor{palette tertiary}{bg=TecAzulOscuro!80!black,fg=white}
\setbeamercolor{title}{fg=TecAzulOscuro,bg=TecFondoBloque}
\setbeamercolor{frametitle}{fg=TecAzulOscuro,bg=TecFondoBloque}

% Bloques personalizados elegantes
\setbeamercolor{block title}{bg=TecAzulOscuro,fg=white}
\setbeamercolor{block body}{bg=TecFondoBloque,fg=black}
\setbeamercolor{block title alerted}{bg=TecRojo,fg=white}
\setbeamercolor{block body alerted}{bg=TecRojo!8,fg=black}
\setbeamercolor{block title example}{bg=TecAzul,fg=white}
\setbeamercolor{block body example}{bg=TecAzul!8,fg=black}

% Redondear bordes y añadir sombra sutil
\setbeamertemplate{blocks}[rounded][shadow=true]
\setbeamertemplate{navigation symbols}{}

% Pie de página limpio con número de diapositiva
\setbeamertemplate{footline}{
  \leavevmode%
  \hbox{%
  \begin{beamercolorbox}[wd=.7\paperwidth,ht=2.25ex,dp=1ex,left]{author in head/foot}%
    \hspace*{2ex}\insertshortauthor~~(\insertshortinstitute) \hfill \insertshorttitle
  \end{beamercolorbox}%
  \begin{beamercolorbox}[wd=.3\paperwidth,ht=2.25ex,dp=1ex,right]{date in head/foot}%
    \insertshortdate{}\hspace*{2em}
    \insertframenumber{} / \inserttotalframenumber\hspace*{2ex}
  \end{beamercolorbox}}%
  \vskip0pt%
}

% Configuración de Listings de Python para visualización pedagógica de código
\lstset{
    language=Python,
    basicstyle=\ttfamily\footnotesize,
    keywordstyle=\color{TecAzulOscuro}\bfseries,
    stringstyle=\color{TecRojo},
    commentstyle=\color{TecGris}\itshape,
    showstringspaces=false,
    breaklines=true,
    frame=lines,
    rulecolor=\color{TecAzulOscuro!30},
    backgroundcolor=\color{TecFondoBloque},
    aboveskip=0.5em,
    belowskip=0.5em,
    literate={á}{{\'a}}1 {é}{{\'e}}1 {í}{{\'i}}1 {ó}{{\'o}}1 {ú}{{\'u}}1
             {Á}{{\'A}}1 {É}{{\'E}}1 {Í}{{\'I}}1 {Ó}{{\'O}}1 {Ú}{{\'U}}1
             {ñ}{{\~n}}1 {Ñ}{{\~N}}1 {¿}{{?`}}1 {¡}{{!`}}1
}

% Metadatos institucionales por defecto
\title[Modelación Estadística]{Modelación Estadística para Ciencia de Datos e Ingeniería}
\author[J. Castillo Colmenares]{Juliho David Castillo Colmenares}
\institute[Tec de Monterrey]{Tecnológico de Monterrey \\ Escuela de Ingeniería y Ciencias}
\date{\today}
% Comandos y macros estadísticas compartidas para las presentaciones Beamer (Metropolis)

% Conjuntos numéricos básicos
\newcommand{\R}{\mathbb{R}}
\newcommand{\N}{\mathbb{N}}
\newcommand{\Q}{\mathbb{Q}}
\newcommand{\Z}{\mathbb{Z}}
\newcommand{\set}[1]{\left\{ #1 \right\}}
\newcommand{\abs}[1]{\left|#1\right|}
\newcommand{\norm}[1]{\left\| #1 \right\|}

% Comandos estadísticos y probabilísticos del curso
\newcommand{\Var}{\operatorname{Var}}
\newcommand{\cov}{\operatorname{Cov}}
\newcommand{\corr}{\rho}
\newcommand{\comb}[2]{\begin{pmatrix} #1 \\ #2 \end{pmatrix}}
\newcommand{\s}{\sigma}
\newcommand{\particion}{\mathfrak{P}}
\newcommand{\card}[1]{\#\left( #1 \right)}
\newcommand{\seq}[3]{\set{#1}_{#2}^{#3}}
\newcommand{\E}{\mathbb{E}}
\newcommand{\Prob}{\mathbb{P}}

% Entornos pedagógicos y cajas en español académico natural
\newenvironment{discusion}{%
  \begin{alertblock}{Pregunta de Discusión}%
}{%
  \end{alertblock}%
}

\newenvironment{reflexion}{%
  \begin{alertblock}{Reflexión para el Aula}%
}{%
  \end{alertblock}%
}

\newenvironment{paradoja}{%
  \begin{alertblock}{Paradoja de Exploración}%
}{%
  \end{alertblock}%
}

\newenvironment{motivacion}{%
  \begin{exampleblock}{Motivación: Ciencia de Datos en la Práctica}%
}{%
  \end{exampleblock}%
}

\newenvironment{retoclase}{%
  \begin{block}{Reto Rápido en Equipo}%
}{%
  \end{block}%
}

\title[Estrategias de experimentación]{Sección 08.01: Estrategias de experimentación}\subtitle{Diseño de Experimentos y ANOVA}
\author[J. Castillo Colmenares]{Juliho Castillo Colmenares}\institute{Tecnológico de Monterrey}\date{}
\begin{document}
\begin{frame}[plain]\titlepage\end{frame}
\begin{frame}{Hoja de ruta}\small\begin{enumerate}\item Motivación y terminología.\item Principios de Fisher.\item Planificación y análisis.\end{enumerate}\end{frame}
\begin{frame}{Fuente y alcance}\small Esta presentación sigue la sección de estrategias de experimentación y presenta el puente entre diseño de Experimentos (DoE) y ANOVA.\begin{block}{Pregunta guía}¿Cómo diseñar datos que permitan distinguir efectos causales de ruido?\end{block}\end{frame}
\begin{frame}{Por qué experimentar}\small Un estudio observacional puede mezclar tratamientos con variables de confusión. La asignación controlada permite comparar escenarios bajo un protocolo reproducible.\end{frame}
\begin{frame}{Inflación por comparaciones}\small Comparar $k$ grupos con pruebas por pares multiplica oportunidades de falso positivo. Para $m$ comparaciones independientes,\[\alpha_{\mathrm{global}}=1-(1-\alpha)^m.\]\end{frame}
\begin{frame}{Terminología esencial}\small\begin{itemize}\item Unidad experimental: entidad que recibe un tratamiento.\item Factor: variable manipulada.\item Tratamiento: combinación concreta de niveles.\item Respuesta: variable medida.\end{itemize}\end{frame}
\begin{frame}{Replicación}\small Repetir un tratamiento en unidades independientes estima la varianza del error y reduce el error estándar de la media. Sin replicación no se puede separar efecto y ruido.\end{frame}
\begin{frame}{Aleatorización}\small Asignar tratamientos al azar evita que diferencias sistemáticas de las unidades se confundan con efectos del factor. También justifica la referencia probabilística del análisis.\end{frame}
\begin{frame}{Bloqueo}\small Agrupar unidades relativamente homogéneas antes de asignar tratamientos controla una fuente conocida de variabilidad y deja un error residual más pequeño.\end{frame}
\begin{frame}{Protocolo experimental}\small\begin{enumerate}\item Definir pregunta y respuesta.\item Elegir factores y niveles.\item Fijar unidades, replicación y bloques.\item Aleatorizar la asignación.\end{enumerate}\end{frame}
\begin{frame}{Puente hacia ANOVA}\small Un diseño bien planificado produce medias comparables y una estimación interna del error. ANOVA contrasta varios tratamientos en un único modelo.\end{frame}
\begin{frame}{Matriz de diseño}\small La matriz registra cada unidad, sus niveles de factor, el bloque si existe, el orden aleatorio y la respuesta observada. Es la base de la trazabilidad.\end{frame}
\begin{frame}{Aplicación: prueba A/B/n}\small Se desean comparar varias configuraciones de un algoritmo. Cada configuración es un tratamiento; el tiempo de ejecución es la respuesta.\end{frame}
\begin{frame}{Aplicación: riesgo familiar}\small Con $k=5$ tratamientos hay $m=10$ pares. Si cada comparación usa $\alpha=0.05$,\[1-(0.95)^{10}\approx0.4013.\]El riesgo global supera al nivel nominal.\end{frame}
\begin{frame}{Aplicación: replicación}\small Ejecutar cada configuración en varias unidades independientes permite estimar la variabilidad del tiempo y distinguir una diferencia estable de una fluctuación.\end{frame}
\begin{frame}{Aplicación: aleatorización}\small Alternar aleatoriamente el orden de ejecución evita que calentamiento, hora o carga del servidor se confundan con el tratamiento.\end{frame}
\begin{frame}{Aplicación: bloqueo}\small Si los servidores pertenecen a generaciones distintas, cada generación puede ser un bloque y las configuraciones se asignan dentro de cada bloque.\end{frame}
\begin{frame}{Aplicación: análisis}\small El plan final combina tratamiento, replicación, aleatorización y bloqueo; después se aplica ANOVA y, si procede, comparaciones ajustadas.\end{frame}
\begin{frame}{Interpretación}\small La inferencia causal depende tanto del diseño como del cálculo. Un valor $p$ no repara una asignación sesgada ni una respuesta mal definida.\end{frame}
\begin{frame}{Comunicación}\small Reporte unidades, factores, niveles, número de réplicas, bloques, regla de aleatorización y respuesta. Así el análisis puede reproducirse.\end{frame}
\begin{frame}{Síntesis}\small\begin{itemize}\item Replicar estima error.\item Aleatorizar protege contra confusión.\item Bloquear controla heterogeneidad conocida.\item ANOVA compara varios tratamientos.\end{itemize}\end{frame}
\begin{frame}{Conexión con el capítulo}\small Las secciones siguientes desarrollan ANOVA de un factor, efectos fijos, diseños bloqueados y diseños factoriales.\end{frame}
\end{document}
